\documentclass[11pt]{article}
\usepackage[margin=1in]{geometry}
\usepackage{amsmath,amsthm,amssymb,mathtools}
\usepackage[T1]{fontenc}
\usepackage{lmodern}
\usepackage{microtype}
\usepackage[hidelinks]{hyperref}
\usepackage{natbib}

\numberwithin{equation}{section}

\newcommand{\E}{\mathbb E}
\newcommand{\PP}{\mathbb P}
\newcommand{\1}{\mathbf 1}
\newcommand{\indep}{\perp\!\!\!\perp}
\DeclareMathOperator{\Cov}{Cov}
\DeclareMathOperator*{\argmin}{arg\,min}

\newtheorem{assumption}{Assumption}[section]
\newtheorem{theorem}[assumption]{Theorem}
\newtheorem{corollary}[assumption]{Corollary}
\newtheorem{proposition}[assumption]{Proposition}
\theoremstyle{remark}
\newtheorem{remark}[assumption]{Remark}

\title{Binary-Instrument IV Poisson Regression}
\author{Lihua Lei}
\date{\today}

\begin{document}
\maketitle

\section{Binary instrument without covariates}

Let $(Y,D,Z)$ denote a generic observation, where $D \in \{0,1\}$,
$Z \in \{0,1\}$, and
\[
Y = Y(D), \qquad D = (1-Z)D(0) + ZD(1), \qquad Y(0),Y(1) \ge 0.
\]
Throughout, $Y(0)$ and $Y(1)$ denote potential outcomes and $D(0)$ and $D(1)$
denote potential treatment states.

\begin{assumption}[No-covariate IV setup]\label{ass:nocov}
The data satisfy:
\begin{itemize}
    \item[(i)] \textbf{Exogeneity:}
    \[
    (Y(0),Y(1),D(0),D(1)) \indep Z.
    \]
    \item[(ii)] \textbf{Relevance:}
    \[
    \E[D\mid Z=1] \neq \E[D\mid Z=0].
    \]
    \item[(iii)] \textbf{Integrability:}
    \[
    \E[Y(0)+Y(1)] < \infty.
    \]
\end{itemize}
\end{assumption}

Fix $\theta \in (-1,\infty)$. The Poisson log-link mean with a binary treatment can be
written as
\[
\exp\{h+\tau D\} = \mu_0(1+\theta D), \qquad \theta = e^{\tau}-1,
\]
where $\mu_0>0$ is the baseline mean. With a binary instrument there is no loss from
using $Z$ itself as the instrument, because the saturated first stage $\E[D\mid Z]$ is an
affine transform of $Z$.

The population Poisson-IV moment system is therefore
\begin{align}
\E\bigl[Y-\mu_0^\star(1+\theta^\star D)\bigr] &= 0, \label{eq:nocov-mean}\\
\E\bigl[Z\{Y-\mu_0^\star(1+\theta^\star D)\}\bigr] &= 0, \label{eq:nocov-iv}
\end{align}
for a pair $(\mu_0^\star,\theta^\star)$ with $\theta^\star \in (-1,\infty)$.

\begin{theorem}\label{thm:nocov-wald}
Suppose Assumption~\ref{ass:nocov} holds and $(\mu_0^\star,\theta^\star)$ satisfies
\eqref{eq:nocov-mean}--\eqref{eq:nocov-iv}. Then
\[
\mu_0^\star = \frac{\E[Y]}{1+\theta^\star \E[D]} \ge 0.
\]
If $\E[Y]>0$, then $\mu_0^\star>0$. In addition,
\begin{equation}\label{eq:nocov-beta}
\mu_0^\star\theta^\star
= \frac{\E[Y\mid Z=1]-\E[Y\mid Z=0]}{\E[D\mid Z=1]-\E[D\mid Z=0]}.
\end{equation}
If $\mu_0^\star>0$, then
\begin{equation}\label{eq:nocov-theta}
\theta^\star
= \frac{\beta_{\mathrm{IV}}}{\E[Y]-\E[D]\beta_{\mathrm{IV}}},
\qquad
\beta_{\mathrm{IV}}
:= \frac{\E[Y\mid Z=1]-\E[Y\mid Z=0]}{\E[D\mid Z=1]-\E[D\mid Z=0]}.
\end{equation}
\end{theorem}

\begin{proof}
The first moment \eqref{eq:nocov-mean} gives
\[
\E[Y] = \mu_0^\star\bigl(1+\theta^\star\E[D]\bigr).
\]
Because $\theta^\star>-1$ and $\E[D]\in[0,1]$, the denominator
$1+\theta^\star\E[D]$ is strictly positive. Hence
\[
\mu_0^\star = \frac{\E[Y]}{1+\theta^\star\E[D]} \ge 0.
\]
If $\E[Y]>0$, then $\mu_0^\star>0$ as well.

Now let $\bar z := \E[Z]$. Subtracting $\bar z$ times \eqref{eq:nocov-mean} from
\eqref{eq:nocov-iv} yields
\[
0 = \E\bigl[(Z-\bar z)\{Y-\mu_0^\star(1+\theta^\star D)\}\bigr].
\]
Since $\E[Z-\bar z]=0$, the constant part drops out, so
\[
\Cov(Z,Y) = \mu_0^\star\theta^\star\Cov(Z,D).
\]
By relevance, $\Cov(Z,D)\neq 0$. Because $Z$ is binary,
\[
\Cov(Z,V) = \PP(Z=1)\PP(Z=0)\bigl(\E[V\mid Z=1]-\E[V\mid Z=0]\bigr)
\]
for any integrable random variable $V$. Applying this identity with $V=Y$ and $V=D$
gives \eqref{eq:nocov-beta}.

If $\mu_0^\star>0$, then \eqref{eq:nocov-beta} implies
\[
\mu_0^\star\theta^\star = \beta_{\mathrm{IV}}.
\]
Substituting this into the first moment gives
\[
\E[Y] = \mu_0^\star + \E[D]\beta_{\mathrm{IV}},
\]
so
\[
\mu_0^\star = \E[Y]-\E[D]\beta_{\mathrm{IV}}.
\]
Dividing $\beta_{\mathrm{IV}}=\mu_0^\star\theta^\star$ by this expression proves
\eqref{eq:nocov-theta}.
\end{proof}

Theorem~\ref{thm:nocov-wald} shows that the Poisson-IV coefficient is not a new causal
object. The causal object is the Wald coefficient $\beta_{\mathrm{IV}}$. The Poisson
coefficient $\theta^\star$ simply rescales it by the positive baseline mean $\mu_0^\star$.

\begin{assumption}[Monotonicity]\label{ass:mono-nocov}
\[
D(1) \ge D(0) \quad \text{almost surely.}
\]
\end{assumption}

\begin{corollary}\label{cor:nocov-late}
Suppose Assumptions~\ref{ass:nocov} and \ref{ass:mono-nocov} hold. Let
\[
C := \{D(1)>D(0)\}
\]
denote the complier event. Then
\begin{equation}\label{eq:nocov-late}
\mu_0^\star\theta^\star = \E\bigl[Y(1)-Y(0)\mid C\bigr].
\end{equation}
If $\mu_0^\star>0$, then
\begin{equation}\label{eq:nocov-late-theta}
\theta^\star
= \frac{\E\bigl[Y(1)-Y(0)\mid C\bigr]}{\mu_0^\star}.
\end{equation}
\end{corollary}

\begin{proof}
By exogeneity and the definitions of the observed outcome and treatment,
\[
\E[Y\mid Z=1]-\E[Y\mid Z=0]
= \E\bigl[Y(D(1))-Y(D(0))\bigr].
\]
Because $D$ is binary,
\[
Y(D(1))-Y(D(0)) = \bigl(D(1)-D(0)\bigr)\bigl(Y(1)-Y(0)\bigr).
\]
Under monotonicity, $D(1)-D(0)=\1\{C\}$, so
\[
\E[Y\mid Z=1]-\E[Y\mid Z=0]
= \PP(C)\E\bigl[Y(1)-Y(0)\mid C\bigr].
\]
Similarly,
\[
\E[D\mid Z=1]-\E[D\mid Z=0]
= \E[D(1)-D(0)] = \PP(C).
\]
Taking the ratio shows that $\beta_{\mathrm{IV}}=\E[Y(1)-Y(0)\mid C]$. The claim then
follows from Theorem~\ref{thm:nocov-wald}.
\end{proof}

\begin{remark}\label{rem:nocov-sign}
If $\mu_0^\star>0$, then $\theta^\star$ and the LATE in
\eqref{eq:nocov-late-theta} have the same sign. Thus, in the binary-instrument model
without covariates, Poisson IV preserves the sign of the usual IV estimand.
\end{remark}

\section{Binary instrument with covariates}

Now restore covariates $X$. Write
\[
e(X) := \PP(Z=1\mid X),
\]
and, for $z\in\{0,1\}$,
\[
m_z(X) := \E[Y\mid Z=z,X], \qquad p_z(X) := \E[D\mid Z=z,X].
\]

\begin{assumption}[Covariate IV setup]\label{ass:cov}
The data satisfy:
\begin{itemize}
    \item[(i)] \textbf{Conditional exogeneity:}
    \[
    (Y(0),Y(1),D(0),D(1)) \indep Z \mid X.
    \]
    \item[(ii)] \textbf{Overlap:}
    \[
    e(X) \in (0,1) \quad \text{almost surely.}
    \]
    \item[(iii)] \textbf{Integrability:}
    \[
    \E[Y(0)+Y(1)] < \infty.
    \]
\end{itemize}
\end{assumption}

\begin{assumption}[Conditional monotonicity]\label{ass:mono-cov}
\[
D(1) \ge D(0) \quad \text{almost surely.}
\]
\end{assumption}

\subsection{Saturated covariates}

The clean benchmark lets the baseline mean be an unrestricted function of $X$.
Let $(\mu_0^\star(X),\theta^\star)$ satisfy
\begin{align}
\E\bigl[Y-\mu_0^\star(X)(1+\theta^\star D)\mid X\bigr] &= 0
\quad \text{almost surely}, \label{eq:cov-mean}\\
\E\bigl[(Z-e(X))\{Y-\mu_0^\star(X)(1+\theta^\star D)\}\bigr] &= 0.
\label{eq:cov-iv}
\end{align}
The first equation is the unrestricted control-function condition. The second is the
usual binary-instrument moment after residualizing the instrument nonparametrically.

\begin{theorem}\label{thm:cov-saturated}
Suppose Assumption~\ref{ass:cov} holds and $(\mu_0^\star(X),\theta^\star)$ satisfies
\eqref{eq:cov-mean}--\eqref{eq:cov-iv}. Then
\[
\mu_0^\star(X)
= \frac{\E[Y\mid X]}{1+\theta^\star\E[D\mid X]} \ge 0
\quad \text{almost surely}.
\]
In addition,
\begin{equation}\label{eq:cov-basic}
\E\Bigl[e(X)(1-e(X))\bigl\{m_1(X)-m_0(X)-\mu_0^\star(X)\theta^\star
\bigl(p_1(X)-p_0(X)\bigr)\bigr\}\Bigr] = 0.
\end{equation}
Define
\[
\beta(X) :=
\begin{cases}
\dfrac{m_1(X)-m_0(X)}{p_1(X)-p_0(X)}, & p_1(X)\neq p_0(X),\\[1.25ex]
0, & p_1(X)=p_0(X).
\end{cases}
\]
If
\begin{equation}\label{eq:cov-positive}
\E\Bigl[e(X)(1-e(X))\mu_0^\star(X)\bigl(p_1(X)-p_0(X)\bigr)\Bigr] > 0,
\end{equation}
then
\begin{equation}\label{eq:cov-ratio}
\theta^\star
= \frac{\E\Bigl[e(X)(1-e(X))\bigl(p_1(X)-p_0(X)\bigr)\beta(X)\Bigr]}
{\E\Bigl[e(X)(1-e(X))\mu_0^\star(X)\bigl(p_1(X)-p_0(X)\bigr)\Bigr]}.
\end{equation}
\end{theorem}

\begin{proof}
The conditional mean moment \eqref{eq:cov-mean} implies
\[
\E[Y\mid X] = \mu_0^\star(X)\bigl(1+\theta^\star\E[D\mid X]\bigr)
\quad \text{almost surely}.
\]
Because $\theta^\star>-1$ and $\E[D\mid X]\in[0,1]$, the multiplier
$1+\theta^\star\E[D\mid X]$ is strictly positive almost surely. Hence
\[
\mu_0^\star(X)
= \frac{\E[Y\mid X]}{1+\theta^\star\E[D\mid X]} \ge 0
\quad \text{almost surely}.
\]

Next condition on $X$ in \eqref{eq:cov-iv}. Since $\mu_0^\star(X)$ is measurable with
respect to $X$ and $\E[Z-e(X)\mid X]=0$,
\[
\E\bigl[(Z-e(X))Y\mid X\bigr]
= \mu_0^\star(X)\theta^\star\E\bigl[(Z-e(X))D\mid X\bigr].
\]
Because $Z$ is binary,
\[
\E\bigl[(Z-e(X))Y\mid X\bigr] = e(X)(1-e(X))\bigl(m_1(X)-m_0(X)\bigr)
\]
and
\[
\E\bigl[(Z-e(X))D\mid X\bigr] = e(X)(1-e(X))\bigl(p_1(X)-p_0(X)\bigr).
\]
Substituting these identities and taking expectations proves \eqref{eq:cov-basic}.

Finally, by the definition of $\beta(X)$,
\[
\bigl(p_1(X)-p_0(X)\bigr)\beta(X) = m_1(X)-m_0(X)
\quad \text{almost surely}.
\]
Replacing $m_1(X)-m_0(X)$ by this product in \eqref{eq:cov-basic} and rearranging
establishes \eqref{eq:cov-ratio} whenever the denominator in
\eqref{eq:cov-positive} is positive.
\end{proof}

Theorem~\ref{thm:cov-saturated} is the direct analogue of the no-covariate result:
$\theta^\star$ is obtained by averaging conditional Wald coefficients and then dividing
by a positive baseline mean.

\begin{corollary}\label{cor:cov-late}
Suppose Assumptions~\ref{ass:cov} and \ref{ass:mono-cov} hold, and assume
\eqref{eq:cov-positive}. Let
\[
C := \{D(1)>D(0)\}.
\]
Then
\[
p_1(X)-p_0(X) = \PP(C\mid X)
\quad \text{almost surely},
\]
and one may take
\[
\beta(X) =
\begin{cases}
\E\bigl[Y(1)-Y(0)\mid C,X\bigr], & \PP(C\mid X)>0,\\[0.75ex]
0, & \PP(C\mid X)=0.
\end{cases}
\]
Hence
\begin{equation}\label{eq:cov-late-ratio}
\theta^\star
= \frac{\E\Bigl[e(X)(1-e(X))\PP(C\mid X)\beta(X)\Bigr]}
{\E\Bigl[e(X)(1-e(X))\PP(C\mid X)\mu_0^\star(X)\Bigr]}.
\end{equation}
Equivalently,
\begin{equation}\label{eq:cov-late-avg}
\theta^\star = \E\bigl[\omega^\star(X)\vartheta(X)\bigr],
\end{equation}
where
\[
\vartheta(X) := \frac{\beta(X)}{\mu_0^\star(X)}
\]
on the event $\{\mu_0^\star(X)>0\}$, and
\[
\omega^\star(X)
:= \frac{e(X)(1-e(X))\PP(C\mid X)\mu_0^\star(X)}
{\E\Bigl[e(X)(1-e(X))\PP(C\mid X)\mu_0^\star(X)\Bigr]}.
\]
The weight $\omega^\star(X)$ is nonnegative and integrates to one.
\end{corollary}

\begin{proof}
By conditional exogeneity,
\[
m_1(X)-m_0(X)
= \E\bigl[Y(D(1))-Y(D(0))\mid X\bigr].
\]
As in the no-covariate case,
\[
Y(D(1))-Y(D(0)) = \bigl(D(1)-D(0)\bigr)\bigl(Y(1)-Y(0)\bigr).
\]
Under monotonicity, $D(1)-D(0)=\1\{C\}$, so
\[
m_1(X)-m_0(X)
= \PP(C\mid X)\E\bigl[Y(1)-Y(0)\mid C,X\bigr]
\]
on the event $\{\PP(C\mid X)>0\}$, and $m_1(X)-m_0(X)=0$ when $\PP(C\mid X)=0$.
Similarly,
\[
p_1(X)-p_0(X) = \E[D(1)-D(0)\mid X] = \PP(C\mid X).
\]
Thus the conditional Wald coefficient $\beta(X)$ equals the conditional LATE on the
cells where the first stage is nonzero, and equals zero elsewhere by convention.
Substituting these identities into \eqref{eq:cov-ratio} gives
\eqref{eq:cov-late-ratio}. Writing the ratio as a normalized weighted average proves
\eqref{eq:cov-late-avg}.
\end{proof}

\begin{remark}\label{rem:cov-benchmark}
When $X$ is degenerate, Theorem~\ref{thm:cov-saturated} and
Corollary~\ref{cor:cov-late} collapse to Theorem~\ref{thm:nocov-wald} and
Corollary~\ref{cor:nocov-late}. The only difference introduced by covariates is that
one averages the conditional Wald or conditional LATE objects with weights
proportional to $e(X)(1-e(X))\mu_0^\star(X)\PP(C\mid X)$.
\end{remark}

\subsection{Linear controls and rich covariates}

Now let $R(X)$ be a vector of included controls containing a constant. Consider the
Poisson-IV specification
\[
\mu_0^\star(X) := \exp\{R(X)'\alpha^\star\},
\]
where $(\alpha^\star,\theta^\star)$ satisfies the population first-order conditions
\begin{align}
\E\Bigl[R(X)\bigl\{Y-\mu_0^\star(X)(1+\theta^\star D)\bigr\}\Bigr] &= 0,
\label{eq:lin-mean}\\
\E\Bigl[Z\bigl\{Y-\mu_0^\star(X)(1+\theta^\star D)\bigr\}\Bigr] &= 0.
\label{eq:lin-iv}
\end{align}
Because $\mu_0^\star(X)$ is exponential, it is strictly positive automatically.

Let $L(X)$ denote the linear projection of $Z$ onto $R(X)$:
\[
L(X) := R(X)'\lambda^\star,
\qquad
\lambda^\star := \argmin_{\lambda}\E\bigl[(Z-R(X)'\lambda)^2\bigr].
\]
Equivalently,
\[
\E\bigl[R(X)\{Z-L(X)\}\bigr] = 0.
\]
Following \cite{blandhol2022tsls}, call the condition
\begin{equation}\label{eq:rich}
e(X) = L(X) \quad \text{almost surely}
\end{equation}
\emph{rich covariates}.

\begin{theorem}\label{thm:rich-covariates}
Suppose Assumption~\ref{ass:cov} holds, $(\alpha^\star,\theta^\star)$ satisfies
\eqref{eq:lin-mean}--\eqref{eq:lin-iv}, and the rich-covariates condition
\eqref{eq:rich} holds. Then
\begin{equation}\label{eq:rich-basic}
\E\Bigl[e(X)(1-e(X))\bigl\{m_1(X)-m_0(X)-\mu_0^\star(X)\theta^\star
\bigl(p_1(X)-p_0(X)\bigr)\bigr\}\Bigr] = 0.
\end{equation}
If
\[
\E\Bigl[e(X)(1-e(X))\mu_0^\star(X)\bigl(p_1(X)-p_0(X)\bigr)\Bigr] > 0,
\]
then
\begin{equation}\label{eq:rich-ratio}
\theta^\star
= \frac{\E\Bigl[e(X)(1-e(X))\bigl(p_1(X)-p_0(X)\bigr)\beta(X)\Bigr]}
{\E\Bigl[e(X)(1-e(X))\mu_0^\star(X)\bigl(p_1(X)-p_0(X)\bigr)\Bigr]},
\end{equation}
where
\[
\beta(X) :=
\begin{cases}
\dfrac{m_1(X)-m_0(X)}{p_1(X)-p_0(X)}, & p_1(X)\neq p_0(X),\\[1.25ex]
0, & p_1(X)=p_0(X).
\end{cases}
\]
Under Assumptions~\ref{ass:cov} and \ref{ass:mono-cov}, the same formula holds with
$\beta(X)$ equal to the conditional LATE on the cells where $\PP(C\mid X)>0$.
\end{theorem}

\begin{proof}
Let
\[
U^\star := Y-\mu_0^\star(X)(1+\theta^\star D).
\]
Because $L(X)$ lies in the linear span of $R(X)$, the first-order condition
\eqref{eq:lin-mean} implies
\[
\E\bigl[L(X)U^\star\bigr] = 0.
\]
Subtracting this equation from \eqref{eq:lin-iv} gives
\[
\E\bigl[(Z-L(X))U^\star\bigr] = 0.
\]
Under rich covariates, $L(X)=e(X)$ almost surely, so
\[
\E\bigl[(Z-e(X))U^\star\bigr] = 0.
\]
Since $\mu_0^\star(X)$ is measurable with respect to $X$ and
$\E[Z-e(X)\mid X]=0$, the same conditional calculation as in the proof of
Theorem~\ref{thm:cov-saturated} yields
\eqref{eq:rich-basic}. Rearranging proves \eqref{eq:rich-ratio}. The final claim then
follows from the conditional Wald-LATE identity used in the proof of
Corollary~\ref{cor:cov-late}.
\end{proof}

Theorem~\ref{thm:rich-covariates} is the Poisson analogue of the familiar saturated or
rich-covariates result for linear IV: once the instrument is residualized correctly, the
coefficient is built from conditional Wald effects alone.

\begin{proposition}\label{prop:level-dependence}
Suppose Assumption~\ref{ass:cov} holds and $(\alpha^\star,\theta^\star)$ satisfies
\eqref{eq:lin-mean}--\eqref{eq:lin-iv}. Then, without imposing rich covariates,
\begin{align}
0
&= \E\Bigl[e(X)(1-e(X))\bigl\{m_1(X)-m_0(X)-\mu_0^\star(X)\theta^\star
\bigl(p_1(X)-p_0(X)\bigr)\bigr\}\Bigr] \notag\\
&\qquad
+ \E\Bigl[(e(X)-L(X))\bigl\{\E[Y\mid X]-\mu_0^\star(X)(1+\theta^\star\E[D\mid X])\bigr\}\Bigr].
\label{eq:level-dependence}
\end{align}
The second term is a level term. It vanishes automatically under rich covariates, but in
general it depends on the levels $\E[Y\mid X]$ and $\E[D\mid X]$, not just on the
conditional Wald contrasts $m_1(X)-m_0(X)$ and $p_1(X)-p_0(X)$.
\end{proposition}

\begin{proof}
As in the proof of Theorem~\ref{thm:rich-covariates}, let
$U^\star=Y-\mu_0^\star(X)(1+\theta^\star D)$. Since
$\E[L(X)U^\star]=0$ and $\E[ZU^\star]=0$,
\[
\E\bigl[(Z-L(X))U^\star\bigr] = 0.
\]
Conditioning on $X$ gives
\[
\E\bigl[(Z-L(X))U^\star\mid X\bigr]
= \Cov(Z,U^\star\mid X) + (e(X)-L(X))\E[U^\star\mid X].
\]
Because $\mu_0^\star(X)$ is measurable with respect to $X$,
\[
\Cov(Z,U^\star\mid X)
= \Cov(Z,Y\mid X) - \mu_0^\star(X)\theta^\star\Cov(Z,D\mid X).
\]
For binary $Z$,
\[
\Cov(Z,Y\mid X) = e(X)(1-e(X))\bigl(m_1(X)-m_0(X)\bigr)
\]
and
\[
\Cov(Z,D\mid X) = e(X)(1-e(X))\bigl(p_1(X)-p_0(X)\bigr).
\]
Moreover,
\[
\E[U^\star\mid X]
= \E[Y\mid X]-\mu_0^\star(X)(1+\theta^\star\E[D\mid X]).
\]
Substituting these expressions into the conditional identity above and taking
expectations yields \eqref{eq:level-dependence}.
\end{proof}

\begin{remark}\label{rem:rich-examples}
Two familiar cases guarantee rich covariates. First, if $R(X)$ is saturated in $X$,
then $L(X)=\E[Z\mid X]=e(X)$ by construction. Second, if $Z$ is mean independent of
$X$, then $e(X)$ is constant, and any specification containing a constant is rich. These
are exactly the binary-instrument cases emphasized in \cite{blandhol2022tsls}.
\end{remark}

\bibliography{poisson_reg}
\bibliographystyle{plainnat}
\end{document}

%%% Local Variables:
%%% mode: latex
%%% TeX-master: t
%%% End:
